\documentclass[12pt,a4paper]{article}
\usepackage[margin=2.5cm]{geometry}
\usepackage{booktabs}
\usepackage{xcolor}
\usepackage{hyperref}
\usepackage{enumitem}
\usepackage{titlesec}
\usepackage{fancyhdr}
\usepackage{amsmath}

\definecolor{rapidoYellow}{HTML}{FFD700}
\definecolor{darkbg}{HTML}{1A1A2E}
\definecolor{accentBlue}{HTML}{0F3460}

\pagestyle{fancy}
\fancyhf{}
\rhead{DVA Capstone 2 --- SectionA G1}
\lhead{Rapido Ride Analysis}
\cfoot{\thepage}

\title{\textbf{\Huge Rapido Ride Analysis}\\[0.5em]
\large Understanding Demand, Cancellations \& Revenue in Bangalore\\[0.3em]
\normalsize NST DVA Capstone 2 $\cdot$ Section A $\cdot$ Group 1}
\author{Aniket Pathak $\cdot$ Nipun Patlori $\cdot$ Saswataduity Bhuin $\cdot$ Akshay Y\\
Narendra Singh $\cdot$ Sayan Bhattacharya $\cdot$ Vedant Madne}
\date{April 29, 2026}

\begin{document}
\maketitle
\tableofcontents
\newpage

%% ============================================================
\section{Project Overview}
%% ============================================================

\subsection{Business Problem}
Ride-hailing platforms like Rapido face a persistent challenge: balancing rider demand with driver supply across different service types and time windows. High cancellation rates erode customer trust, reduce revenue, and signal inefficiencies in resource allocation.

\textbf{Core Business Question:}\\
\textit{How do service type, time of day, and location influence ride cancellation rates and revenue --- and where should Rapido focus driver allocation to minimise cancellations during peak demand?}

\subsection{Dataset at a Glance}
\begin{center}
\begin{tabular}{ll}
\toprule
\textbf{Attribute} & \textbf{Value} \\
\midrule
Source & Bangalore Rapido Ride Data (Kaggle, simulated) \\
Raw Rows & 51,000 \\
Cleaned Rows & 49,801 \\
Raw Columns & 13 \\
Cleaned Columns & 30 \\
Time Period & April 10 -- April 29, 2024 \\
Format & CSV \\
\bottomrule
\end{tabular}
\end{center}

%% ============================================================
\section{Pipeline: From Raw to Clean (ETL)}
%% ============================================================

The entire pipeline runs across five Jupyter notebooks.

\subsection{Notebook 01 --- Extraction}
\begin{itemize}
  \item Loaded \texttt{rapido\_rides.csv} (51,000 rows, 13 columns) with \texttt{pd.read\_csv()}.
  \item Ran \texttt{df.info()} and \texttt{df.describe()} to inspect dtypes, ranges, and nulls.
  \item Identified the 13 raw columns: \texttt{services, date, time, ride\_status, source, destination, duration, ride\_id, distance, ride\_charge, misc\_charge, total\_fare, payment\_method}.
  \item Initial null audit showed \textbf{10,346 nulls} in \texttt{ride\_charge}, \texttt{misc\_charge}, \texttt{total\_fare} and \textbf{12,398 nulls} in \texttt{payment\_method}.
\end{itemize}

\subsection{Notebook 02 --- Cleaning \& Feature Engineering}

\subsubsection{Step 1: Handling Missing Values (Business-Logic Driven)}
The 10,346 missing fare rows corresponded to \textbf{cancelled rides} --- rides that were booked but never completed, so they naturally have no fare. We did \textbf{NOT} drop them.

\begin{itemize}
  \item Set \texttt{ride\_charge}, \texttt{misc\_charge}, \texttt{total\_fare} $\leftarrow 0$ for cancelled rides.
  \item Set \texttt{payment\_method} $\leftarrow$ \texttt{"not\_applicable"} for cancelled rides.
  \item Recalculated \texttt{total\_fare = ride\_charge + misc\_charge} for completed rides to fix any inconsistencies.
\end{itemize}

\textbf{Key interview point:} We chose business-logic imputation over mean/median imputation because a cancelled ride genuinely has zero fare --- filling with an average would pollute revenue KPIs.

\subsubsection{Step 2: Datetime Conversion}
\begin{itemize}
  \item \texttt{date} column converted from \texttt{object} $\rightarrow$ \texttt{datetime64}.
  \item Combined \texttt{date} + \texttt{time} strings into a new \texttt{timestamp} column for time-series analysis.
\end{itemize}

\subsubsection{Step 3: Duplicate Removal}
\begin{itemize}
  \item Deduplicated on \texttt{ride\_id}, keeping the first occurrence.
  \item Dropped 1,199 duplicate records.
\end{itemize}

\subsubsection{Step 4: Categorical Standardisation}
\begin{itemize}
  \item All text columns lowercased and stripped of whitespace.
  \item Service names normalised: \texttt{"bike lite"} $\rightarrow$ \texttt{bike\_lite}, \texttt{"cab economy"} $\rightarrow$ \texttt{cab\_economy}.
\end{itemize}

\subsubsection{Step 5: Anomaly Filtering}
Dropped completed rides where:
\begin{itemize}
  \item \texttt{distance > 100 km} (physically implausible in Bangalore city)
  \item \texttt{distance $\leq$ 0 km}
  \item \texttt{duration $\leq$ 0 min}
\end{itemize}

\subsubsection{Step 6: Feature Engineering (17 New Columns)}
\begin{center}
\begin{tabular}{lll}
\toprule
\textbf{Column} & \textbf{Formula / Logic} & \textbf{Purpose} \\
\midrule
\texttt{year} & \texttt{date.dt.year} & Temporal filter \\
\texttt{month} & \texttt{date.dt.month} & Seasonal trends \\
\texttt{hour} & \texttt{timestamp.dt.hour} & Hourly demand \\
\texttt{day\_of\_week} & \texttt{timestamp.dt.day\_name()} & Weekly patterns \\
\texttt{is\_weekend} & 1 if Sat/Sun, else 0 & Weekend vs weekday \\
\texttt{peak\_hour} & 1 if hour $\in$ [8,9,10,17,18,19] & Peak flagging \\
\texttt{completed} & 1 if completed, 0 if cancelled & Binary target \\
\texttt{avg\_speed\_kmh} & distance / (duration/60) & Operational efficiency \\
\texttt{fare\_per\_km} & total\_fare / distance & Unit economics \\
\texttt{source\_cancel\_rate} & avg cancel rate per pickup zone & Geographic reliability \\
\texttt{demand\_intensity} & Z-score of volume per source-hour & Supply pressure \\
\texttt{demand\_zone} & low / medium / high demand label & Zone segmentation \\
\texttt{distance\_category} & short / medium / long & Ride length segment \\
\texttt{fare\_category} & budget / standard / premium / cancelled & Revenue segment \\
\bottomrule
\end{tabular}
\end{center}

%% ============================================================
\section{KPI Framework}
%% ============================================================

\begin{center}
\begin{tabular}{llll}
\toprule
\textbf{KPI} & \textbf{Value} & \textbf{Formula} & \textbf{Insight} \\
\midrule
Cancellation Rate & \textbf{10.40\%} & 5,180 / 49,801 & 1 in 10 rides cancelled \\
Total Revenue & \textbf{Rs.68,37,929} & Sum of \texttt{total\_fare} & Completed rides only \\
Avg Revenue/Ride & \textbf{Rs.153.24} & Revenue / completed & Baseline fare benchmark \\
Peak Cancel Rate & \textbf{13.22\%} & 1,653 / 12,503 & 40\% higher than off-peak \\
Off-Peak Cancel Rate & \textbf{9.46\%} & 3,527 / 37,298 & Baseline comparison \\
\bottomrule
\end{tabular}
\end{center}

%% ============================================================
\section{Exploratory Data Analysis (EDA)}
%% ============================================================

\subsection{Service Type Analysis}
\begin{center}
\begin{tabular}{lrrrrr}
\toprule
\textbf{Service} & \textbf{Rides} & \textbf{Cancelled} & \textbf{Cancel \%} & \textbf{Revenue (Rs.)} & \textbf{Avg Fare} \\
\midrule
auto & 11,474 & 1,180 & 10.28\% & 19,36,736 & 188.14 \\
bike & 17,492 & 1,813 & 10.36\% & 13,02,960 & 83.10 \\
bike\_lite & 4,541 & 452 & 9.95\% & 2,32,899 & 56.96 \\
cab\_economy & 9,370 & 1,006 & 10.74\% & 26,98,620 & 322.65 \\
parcel & 6,924 & 729 & 10.53\% & 6,66,713 & 107.62 \\
\bottomrule
\end{tabular}
\end{center}

\textbf{Key insight:} \texttt{cab\_economy} has the highest cancellation rate (10.74\%) and highest average fare (Rs.322.65) --- premium rides are most at risk. \texttt{bike\_lite} is the most reliable service.

\subsection{Peak vs Off-Peak Analysis}
\begin{center}
\begin{tabular}{lrrrr}
\toprule
\textbf{Window} & \textbf{Rides} & \textbf{Cancelled} & \textbf{Cancel Rate} \\
\midrule
Off-Peak & 37,298 & 3,527 & 9.46\% \\
Peak (8-10 AM, 5-7 PM) & 12,503 & 1,653 & \textbf{13.22\%} \\
\bottomrule
\end{tabular}
\end{center}

Peak hours see \textbf{39.7\% more cancellations} than off-peak. The 5 PM -- 7 PM evening slot is the worst at \textbf{14.14\%} cancellation rate.

\subsection{Hour-wise Cancellation Heatmap (Key Hours)}
\begin{center}
\begin{tabular}{lrl}
\toprule
\textbf{Hour} & \textbf{Cancel Rate} & \textbf{Interpretation} \\
\midrule
5 AM & 12.84\% & Late-night stragglers, low driver availability \\
8 AM & 13.67\% & Morning office rush begins \\
10 AM & 13.24\% & Peak tail \\
17 (5 PM) & 14.14\% & \textbf{Worst hour of day} --- evening rush \\
18 (6 PM) & 12.62\% & Evening peak continues \\
12 PM & 7.07\% & Midday lull, best completion \\
14 (2 PM) & 7.00\% & Lowest cancellation of day \\
\bottomrule
\end{tabular}
\end{center}

\subsection{Geographic Analysis}
\begin{itemize}
  \item \textbf{Top Revenue Zone:} Koramangala --- Rs.2,73,504 (4\% of total revenue from one zone).
  \item Koramangala also has 1,957 rides with a 10.17\% cancellation rate.
  \item The zone naming inconsistency (\texttt{koramangala} vs \texttt{kormangala}) was detected during EDA and standardised.
\end{itemize}

\subsection{Demand Zone Distribution}
\begin{center}
\begin{tabular}{lr}
\toprule
\textbf{Demand Zone} & \textbf{Ride Count} \\
\midrule
High Demand & 17,354 \\
Medium Demand & 17,590 \\
Low Demand & 14,857 \\
\bottomrule
\end{tabular}
\end{center}

\subsection{Efficiency Metrics (Completed Rides)}
\begin{center}
\begin{tabular}{lrrrrr}
\toprule
\textbf{Metric} & \textbf{Mean} & \textbf{Median} & \textbf{Min} & \textbf{Max} \\
\midrule
Distance (km) & 6.99 & 5.48 & 0.8 & 45.0 \\
Duration (min) & 24.09 & 19.0 & 5 & 120 \\
Avg Speed (km/h) & 18.31 & 17.69 & 2.07 & 73.1 \\
Fare per km (Rs.) & 20.94 & 19.85 & 0 & 110.18 \\
Total Fare (Rs.) & 137.31 & 93.84 & 0 & 1,301.58 \\
\bottomrule
\end{tabular}
\end{center}

\subsection{Service Efficiency Comparison}
\begin{center}
\begin{tabular}{lrrrr}
\toprule
\textbf{Service} & \textbf{Avg Fare/km} & \textbf{Avg Speed (km/h)} & \textbf{Avg Fare (Rs.)} \\
\midrule
cab\_economy & Rs.33.77 & 19.59 & 322.65 \\
auto & Rs.28.47 & 18.75 & 188.14 \\
parcel & Rs.20.56 & 17.93 & 107.62 \\
bike & Rs.17.76 & 17.89 & 83.10 \\
bike\_lite & Rs.15.05 & 17.46 & 56.96 \\
\bottomrule
\end{tabular}
\end{center}

\subsection{Payment Method Distribution}
\begin{center}
\begin{tabular}{lr}
\toprule
\textbf{Payment Method} & \textbf{Count} \\
\midrule
Paytm & 11,208 \\
QR Scan & 11,186 \\
GPay & 11,147 \\
Amazon Pay & 11,080 \\
Not Applicable (cancelled) & 5,180 \\
\bottomrule
\end{tabular}
\end{center}

\textbf{Note:} Payment methods are nearly uniformly distributed across all four UPI providers --- no single provider dominates.

%% ============================================================
\section{Statistical Analysis}
%% ============================================================

\subsection{Test 1: Chi-Square --- Service Type vs Cancellation}
\begin{itemize}
  \item \textbf{Null Hypothesis:} Service type and cancellation are independent.
  \item \textbf{Result:} $\chi^2 = 2.42$, $p = 0.659$, $df = 4$
  \item \textbf{Conclusion:} \textbf{Fail to reject $H_0$}. No statistically significant difference in cancellation rates across services. All services cancel at approximately the same rate ($\sim$10.4\%).
\end{itemize}

\subsection{Test 2: T-Test --- Peak vs Off-Peak Cancellations}
\begin{itemize}
  \item \textbf{Null Hypothesis:} Peak and off-peak cancellation rates are equal.
  \item \textbf{Result:} $t = -11.95$, $p < 0.0001$ (statistically significant)
  \item \textbf{Conclusion:} \textbf{Reject $H_0$}. Peak hours have a \textbf{significantly higher} cancellation rate (13.22\% vs 9.46\%). The time-of-day effect is real and not due to chance.
\end{itemize}

\subsection{Test 3: Correlation --- Demand Intensity vs Cancellation}
\begin{itemize}
  \item Pearson correlation between \texttt{demand\_intensity} (z-score of volume) and cancellation: $r = -0.0134$
  \item \textbf{Interpretation:} Very weak correlation. High demand alone does not predict cancellation. Other factors (driver supply, time of day) are more important drivers.
\end{itemize}

\subsection{Revenue Leakage Simulation}
\begin{itemize}
  \item Average completed ride fare: Rs.153.24
  \item If cancellations in top revenue zones are reduced by 10\%:
  \item \textbf{Potential monthly recovery} $\approx$ Rs.1,20,000 (extrapolated from 20-day dataset)
\end{itemize}

%% ============================================================
\section{Charts \& Visualisations}
%% ============================================================

All charts were created using \textbf{Python (matplotlib / seaborn)} in \texttt{03\_EDA.ipynb} and \texttt{04\_statistical\_analysis.ipynb}, and published as an interactive \textbf{Tableau Dashboard}.

\subsection{Chart List}
\begin{enumerate}
  \item \textbf{Bar Chart --- Rides by Service Type:} bike leads with 17,492 rides; bike\_lite is smallest with 4,541.
  \item \textbf{Stacked Bar --- Cancellation by Service:} Shows completed vs cancelled split per service. All services cluster near 10\%.
  \item \textbf{Line Chart --- Hourly Cancellation Rate:} Clear spikes at 8 AM (13.67\%) and 5 PM (14.14\%). Troughs at 12--14 PM.
  \item \textbf{Heatmap --- Service $\times$ Peak Hour Cancellation:} 5x2 grid showing peak vs off-peak rates per service. cab\_economy peak = 13.12\%.
  \item \textbf{Box Plot --- Fare Distribution by Service:} cab\_economy has highest spread (Rs.0--1301). bike\_lite is tightest.
  \item \textbf{Scatter Plot --- Distance vs Fare:} Positive correlation, with cab\_economy having steepest fare-per-km slope.
  \item \textbf{Bar Chart --- Revenue by Service:} cab\_economy = Rs.26.9L despite fewer rides; bike = Rs.13L with most rides.
  \item \textbf{Map / Geographic Chart (Tableau):} Cancellation hotspots by pickup zone overlaid on Bangalore map.
  \item \textbf{Histogram --- Fare Distribution:} Right-skewed distribution; most rides between Rs.50--200.
  \item \textbf{Day-of-Week Chart:} Tuesday has highest cancellation rate (11.19\%); Wednesday is most reliable (9.60\%).
\end{enumerate}

\subsection{Tableau Dashboard}
Two dashboard views were published:
\begin{itemize}
  \item \textbf{Executive View:} Cancellation rates by service, peak vs off-peak summary, top revenue zones.
  \item \textbf{Operational View:} Ride volume by hour, cancellation heatmap, fare distribution by service.
\end{itemize}

%% ============================================================
\section{Key Business Insights}
%% ============================================================

\begin{enumerate}
  \item \textbf{Peak Hour is the Primary Risk Factor:} Cancellation spikes 39.7\% during morning (8--10 AM) and evening (5--7 PM) peaks. Time of day, not service type, is the dominant driver.
  \item \textbf{Service Type Doesn't Matter for Cancellation:} Chi-square test confirms all services cancel at the same rate. Operational improvements must focus on time/location, not service segmentation.
  \item \textbf{cab\_economy is the Highest Revenue Risk:} It earns the most per ride (Rs.322.65 avg) and has the highest cancellation rate (10.74\%). Every lost cab ride costs 4x more than a lost bike ride.
  \item \textbf{bike and bike\_lite are the Most Resilient:} During peak hours, bikes maintain relatively low cancellation and high volume. They can serve as overflow capacity.
  \item \textbf{Koramangala is the Critical Zone:} 4\% of total revenue from one zone. This zone must be prioritised for driver incentives during peak hours.
  \item \textbf{Midday (12--3 PM) is the Efficiency Window:} Lowest cancellation rates (7.0--7.5\%), ideal time for driver training or scheduled maintenance.
  \item \textbf{Revenue Concentration:} Top 5 zones contribute disproportionate revenue --- geographic concentration is a risk factor if driver supply isn't maintained there.
  \item \textbf{Payment is Neutral:} No UPI provider dominates, suggesting users have no strong payment preference. Not a lever for intervention.
\end{enumerate}

%% ============================================================
\section{Business Recommendations}
%% ============================================================

\begin{enumerate}
  \item \textbf{Peak-Hour Driver Surge Incentives:} Offer bonus payouts to drivers who are active between 8--10 AM and 5--7 PM. Target Koramangala and HSR Layout first. Expected impact: reduce peak cancellation from 13.22\% to $\sim$10\%.

  \item \textbf{Zone-Based Driver Pre-positioning:} Use demand zone data to pre-position drivers in high-demand zones 15 minutes before peak hours begin. Implement via the driver app as shift recommendations.

  \item \textbf{cab\_economy Supply Boost:} Since cab rides are highest revenue but most vulnerable, create a dedicated cab driver pool for peak hours in top-5 revenue zones.

  \item \textbf{Bike as Peak Overflow:} During extreme demand, surface bike/auto alternatives faster in the UI when cabs are unavailable. Reduces user frustration and captures revenue that would otherwise be lost.

  \item \textbf{Real-Time Cancellation Monitoring:} Deploy a live dashboard (like our Tableau view) for the ops team to track hourly cancellation rates and trigger driver incentives dynamically.
\end{enumerate}

%% ============================================================
\section{Tech Stack \& Team Roles}
%% ============================================================

\begin{center}
\begin{tabular}{ll}
\toprule
\textbf{Tool} & \textbf{Purpose} \\
\midrule
Python & ETL, Cleaning, EDA, Statistical Analysis \\
Pandas & Data manipulation and transformation \\
SciPy & Statistical tests (chi-square, t-test) \\
Matplotlib / Seaborn & Chart generation \\
Tableau Public & Interactive dashboard \\
GitHub & Version control, collaboration \\
Jupyter Notebooks & Development environment \\
\bottomrule
\end{tabular}
\end{center}

\begin{center}
\begin{tabular}{lll}
\toprule
\textbf{Name} & \textbf{Role} & \textbf{Primary Contribution} \\
\midrule
Aniket Pathak & Project Lead & Problem framing, ETL oversight \\
Nipun Patlori & Data Lead & EDA, cleaning pipeline \\
Saswataduity Bhuin & ETL Lead & Tableau dashboard \\
Akshay Y & Analysis Lead & Statistical analysis \\
Narendra Singh & Visualization Lead & Tableau + PPT \\
Sayan Bhattacharya & Strategy Lead & Recommendations \\
Vedant Madne & PPT \& Quality Lead & Presentation quality \\
\bottomrule
\end{tabular}
\end{center}

%% ============================================================
\section{Quick Revision Cheat Sheet}
%% ============================================================

\begin{itemize}[leftmargin=*]
  \item \textbf{Raw data:} 51,000 rows, 13 columns
  \item \textbf{Cleaned data:} 49,801 rows, 30 columns
  \item \textbf{Missing values:} 10,346 in fare columns (cancelled rides) $\rightarrow$ filled with 0
  \item \textbf{Duplicates removed:} 1,199 rows
  \item \textbf{Cancellation Rate:} 10.40\% (5,180 rides)
  \item \textbf{Peak Cancellation:} 13.22\% vs 9.46\% off-peak ($t$-test: $p < 0.0001$)
  \item \textbf{Service test:} Chi-square $p = 0.66$ --- service type does NOT affect cancellation
  \item \textbf{Total Revenue:} Rs.68,37,929
  \item \textbf{Highest revenue service:} cab\_economy (Rs.26.9L)
  \item \textbf{Most rides service:} bike (17,492 rides)
  \item \textbf{Best hour to ride:} 2 PM --- 7.00\% cancellation rate
  \item \textbf{Worst hour to ride:} 5 PM --- 14.14\% cancellation rate
  \item \textbf{Top zone:} Koramangala --- Rs.2,73,504 revenue
  \item \textbf{Peak hours defined:} 8--10 AM and 5--7 PM
  \item \textbf{Avg speed:} 18.31 km/h
  \item \textbf{Avg fare per km:} Rs.20.94
\end{itemize}

\end{document}
